% unidades/12-koto-no.tex
\chapter{\emoji{scroll} \ruby{こと}{} vs \ruby{の}{} — Nominalización}
\unidadheader{12}{名詞化}{こと vs の}{Convertir verbos en sustantivos: abstracto vs concreto}

\begin{tcolorbox}[vocabbox, title={\emoji{pushpin} Vocabulario}]
\begin{tabularx}{\linewidth}{llX}
\toprule
\textbf{Japonés} & \textbf{Español} & \textbf{Tipo} \\ \midrule
\vocabitem{趣味}{しゅみ}{pasatiempo / hobby}{sustantivo}
\vocabitem{驚く}{おどろく}{sorprenderse}{v. G1}
\vocabitem{得意}{とくい}{ser bueno en algo}{adj. na}
\vocabitem{苦手}{にがて}{no ser bueno en algo}{adj. na}
\vocabitem{気づく}{きづく}{darse cuenta}{v. G1}
\vocabitem{感じる}{かんじる}{sentir}{v. G2}
\vocabitem{笑う}{わらう}{reír}{v. G1}
\bottomrule
\end{tabularx}
\end{tcolorbox}

\subsection*{\emoji{speech-balloon} Frases Clave}
\frase{\ruby{泳}{およ}ぐことが\ruby{得意}{とくい}です。}{Soy bueno nadando. (habilidad general)}
\frase{\ruby{約束}{やくそく}を\ruby{守}{まも}ることは\ruby{大切}{たいせつ}です。}{Cumplir las promesas es importante.}
\frase{\ruby{子供}{こども}が\ruby{笑}{わら}うのを\ruby{見}{み}た。}{Vi reír al niño. (evento percibido)}
\frase{\ruby{鳥}{とり}が\ruby{飛}{と}ぶのが\ruby{見}{み}えた。}{Pude ver pájaros volando.}

\subsection*{\emoji{gear} Gramática}
\patron{こと — nominalización abstracta/general}
  {[V-plain] + こと}
  {\ej{\ruby{日本語}{にほんご}を\ruby{話}{はな}すことが\ruby{好}{す}きです。}{Me gusta hablar japonés. (gusto general)}
   \ej{\ruby{毎日}{まいにち}\ruby{運動}{うんどう}することは\ruby{健康}{けんこう}にいい。}{Hacer ejercicio todos los días es bueno para la salud.}
   \ej{\ruby{彼女}{かのじょ}は\ruby{料理}{りょうり}することが\ruby{得意}{とくい}だ。}{Ella es buena cocinando.}}

\patron{の — nominalización concreta / perceptual}
  {[V-plain] + の + が/を (ver, oír, esperar, percibir)}
  {\ej{\ruby{子供}{こども}が\ruby{走}{はし}るのを\ruby{見}{み}た。}{Vi correr a los niños. (percepción directa)}
   \ej{\ruby{誰}{だれ}かが\ruby{歌}{うた}うのが\ruby{聞}{き}こえた。}{Se escuchó a alguien cantar.}
   \ej{\ruby{電車}{でんしゃ}が\ruby{来}{く}るのを\ruby{待}{ま}っている。}{Estoy esperando que llegue el tren.}}

\comparacion{こと (abstracto) vs の (concreto)}
  {こと\\\small abstracto / reglas / hábitos\\\small \ruby{泳}{およ}ぐことが\ruby{好}{す}き\\\textit{Me gusta nadar (en general)}}
  {の\\\small perceptual / evento observado\\\small \ruby{泳}{およ}ぐのを\ruby{見}{み}た\\\textit{Vi (a alguien) nadar}}
  {こと generaliza una acción o la eleva a principio. の la concretiza: alguien la percibe ocurriendo. Con verbos de percepción (見る, 聞く, 感じる) solo se usa の.}

\coloquial{こと → こと (igual) / の más frecuente en casual}{
  の es más natural en habla cotidiana: 行くの好き。食べるの早い。En formal siempre こと es preferido en frases nominales: 〜することが大切です。}

\culturabox{\ruby{日本語}{にほんご}の\ruby{抽象化}{ちゅうしょうか} — Cómo el japonés abstrae acciones}{
  La nominalización es central en japonés. こと transforma acciones en conceptos, permitiendo hablar de principios y reglas. Frases como 〜することが大切 (es importante hacer ~) son pilares del discurso formal japonés, especialmente en educación y ética profesional.
}

\lectura{\ruby{趣味}{しゅみ}について\\[0.4em]
  \ruby{私}{わたし}の\ruby{趣味}{しゅみ}は\ruby{写真}{しゃしん}を\ruby{撮}{と}ることです。\ruby{美}{うつく}しい\ruby{景色}{けしき}を\ruby{撮}{と}ることが\ruby{楽}{たの}しいです。\ruby{先日}{せんじつ}、\ruby{子供}{こども}たちが\ruby{遊}{あそ}ぶのを\ruby{撮}{と}りました。\ruby{笑}{わら}うのを\ruby{見}{み}るのが\ruby{一番}{いちばん}\ruby{好}{す}きです。}
{Sobre mi pasatiempo\\[0.4em]
  Mi pasatiempo es tomar fotografías. Me gusta tomar fotos de paisajes hermosos. El otro día, fotografié a niños jugando. Lo que más me gusta es verlos reír.}

\begin{tcolorbox}[ejerciciobox, title={\emoji{pencil} Ejercicios}]
\begin{enumerate}
  \item ¿こと o の? \ruby{彼}{かれ}が\ruby{泣}{な}く\_\_\_を\ruby{見}{み}た。
  \item ¿こと o の? \ruby{約束}{やくそく}を\ruby{守}{まも}る\_\_\_は\ruby{大切}{たいせつ}だ。
  \item Crea una oración sobre tu hobby usando 〜することが好きです。
\end{enumerate}
\textbf{Respuestas:}
1) の (percepción directa) \quad
2) こと (principio general) \quad
3) (libre)
\end{tcolorbox}

\begin{tcolorbox}[kanjibox, title={\emoji{mahjong} Kanjis}]
\begin{tabularx}{\linewidth}{cllXX}
\toprule
\textbf{Kanji} & \textbf{On} & \textbf{Kun} & \textbf{Significado} & \textbf{Vocab} \\ \midrule
\kanjirow{趣}{しゅ}{\ruby{おもむ}{き}}{sabor/gusto}{\ruby{趣味}{しゅみ} / \ruby{趣旨}{しゅし}}
\kanjirow{味}{み}{\ruby{あじ}{}}{sabor}{\ruby{趣味}{しゅみ} / \ruby{意味}{いみ}}
\kanjirow{驚}{きょう}{\ruby{おどろ}{く}}{sorprenderse}{\ruby{驚}{おどろ}く / \ruby{驚異}{きょうい}}
\kanjirow{感}{かん}{\ruby{かん}{じる}}{sentir}{\ruby{感じる}{かんじる} / \ruby{感情}{かんじょう}}
\bottomrule
\end{tabularx}
\end{tcolorbox}
